% Chapter3_Methodology.tex-ээс 2026-11-27-нд үгчлэн зөөсөн агентын тэмдэглэл (хэвлэгдэхгүй). Шинэ тэмдэглэлийг энд нэмнэ.

%======================================================================
% ЭНЭ ФАЙЛД БАЙГАА БҮХ ЗҮЙЛ БАГШИД ҮЗҮҮЛЭХ НООРОГТ ХЭВЛЭГДЭНЭ (бүтэн ноорогт ч мөн): 3.1–3.5.
% Зөвхөн бүтэн ноорогт орох үнэлгээ, яриа таних аргын сонголт:
% Chapters/full/ch3_evaluation_and_stt_choice.tex (файлын төгсгөлийн \fullinput мөр).
%
% 2026.11.16 (7), багшийн зөвлөгөөгөөр («хэллэг хэт энгийн», «одоо арга зүй хамгийн чухал»):
% найруулгыг Монгол эрдэм шинжилгээний бичвэрийн хэллэгт ойртуулсан. Агуулга, баримт, тоо,
% бүтэц, шошго (label) өөрчлөгдөөгүй. Тулгуурласан уншлага: Монгол хэл дээрх 22 өгүүлэл
% (reports/Монгол эрдэм шинжилгээний бичвэрийн хэллэг.md, papers/36-*). Баримталсан зүйл:
% дунд урттай өгүүлбэр (15–25 үг), «-в»/«-сан» төгсгөл, «уг», «буй», «тус бүр», «-д тулгуурлан»,
% «-ийн дагуу», «тохиолдолд», «болно»; «бид»-гүй; зураг заахдаа «…-д үзүүлэв».
% Хэрэглээний тохиолдлын диаграммыг 1.3-т шилжүүлсэн (багш: эхний хэсэгт).
% Өмнөх хувилбар: thesis/archive-20261116/chapters-before-teacher-feedback-20261116/.
%
% 2026.11.20: 3.6 «Системийг үнэлэх арга» (төлөвлөгөөний W04-ийн ажил) ба «Бүлгийн дүгнэлт» нэмсэн;
% 3.1, 3.4.2-ын хоёр газрын найруулга. Төлөвлөгөөнөөс түрүүлж хийсэн системийн ажлыг энд оруулаагүй.
% Өмнө/дараа: reports/Хэллэгийн туршилт 2026-11-20.md. Өмнөх хувилбар: thesis/archive-20261120/.
%
% 2026.11.16 (3), хэрэглэгчийн шүүмжээр: (1) мэдлэгийн мод бол энэ ажлын RAG-ийн суурь — 3.4-т
% бүтэц, зураг, хайлтад хэрхэн ашиглахыг алхмаар тайлбарласан; (2) «ирээгүй, байхгүй, хоосон» гэж
% бичихгүй — цуглуулах мэдээллийг «Цаашдын ажил»-д «цуглуулна» гэж бичнэ; (3) «сургуулийн албад»
% биш «сургуулийн зүгээс өгсөн мэдээлэл», «сургуулийн ажилтнууд».
% Кодтой тулгасан (11/16): config/knowledge-topics.json (үндэс, 11 хүрээ, 31 дэд сэдэв, хаягаар оноох
% дүрэм), config/knowledge-domains.json (түлхүүр үг, enabled=false), website_retrieval.py
% (predict_domain, тэнцвэл оноо нэмэхгүй), config/referral.json (topic_offices approved=false).
%
% 2026.11.23 (хэрэглэгч: «thesis-ийг эхнээс дуустал дахин хар, үг сонголт, асуултаа хар, зас»): аудит
% papers/45-thesis-language-audit-20261122.mn.md №65–100 — мэдлэгийн модны алхмыг «зохион бүтээсэн,
% үнэлгээгээр идэвхжүүлнэ» гэж үнэн зөв бичсэн (knowledge-domains.json enabled=false); 3.4.2 ба 3.5-ын таван
% алхмыг холбосон; «Одоогоор … цаашид» хэв маягийг авсан; 3.3-ын дөрвөн хүснэгт ↔ Зураг 3.1; 3.6-г 40 асуулт
% (20/10/10) болгосон; бүтэн ноорогт 3.8–3.9 «Бүлгийн дүгнэлт»-ээс өмнө. Өмнөх хувилбар: thesis/archive-20261123/.
% 2026.11.23 (2), орхигдсон аудитын саналаар: мэдээлэл шинэчлэх механизмын өгүүлбэр (№73, төлөвлөгөөнөөс түрүүлсэн
% ажил) → full/ch3_data_update.tex; 3.6-ын «ажил горилогч» → «ажлын байр сонирхогч» (Хүснэгт 1.1, 3.2-той нэг).
%======================================================================

% 2026.11.20: төлөвлөгөөний W04-ийн ажил («Системийн үнэлгээний асуудал: шалгах асуултын багц бэлдэх»).
% 2026.11.23 (хэрэглэгч: «50-ийг 40 болго, бодит хүн юу асуухыг судалж бич»): шинэ ноорог
% data/eval/system-questions-draft-20261122.json (20 хариултай / 10 ресепшний / 10 хариултгүй, бүгд эцсийн
% үнэлгээнийх, хянаагүй); судалгаа papers/44-real-questions-research-20261122.mn.md (локал).
% Хуучин 50: data/eval/system-questions-draft-20261117.json (20/9/9/12, тохируулах 27, эцсийн 23).
% Багцыг сургалтын албатай тулгаж түгжсэний дараа огноог нь энд нэм. Бүтэн ноорогт доорх
% «Системийн үнэлгээ» хэсэг (full/) бодит асуулт цуглуулах өргөн төлөвлөгөөг агуулна — давхцлыг
% бүтэн ноорогийг засахдаа нэгтгэнэ.

% 2026.11.23 (8): хэрэглэгч 40 асуултыг системийг сайжруулахад ашиглахыг зөвшөөрсөн («its is okay even question is
% changed … do it with those 40 questions»). Иймд «Уг багцыг системийг тохируулахад ашиглахгүй» гэсэн амлалтыг
% хасаж, эцсийн үнэлгээг тохируулахад ашиглаагүй бодит асуултаар хийнэ гэж бичив. Энэ 40-ийн дүнг 4-р бүлэгт
% оруулбал «хөгжүүлэлтийн явцад ашигласан багцын дүн» гэж заавал тэмдэглэнэ (runtime/eval/q40-*.json).

%======================================================================
% 2026.11.23 (багш: «Whisper, Gemini хоёрыг 10–15 бичлэг дээр ижил аудиогоор туршиж, үр дүн нь ийм гарсан
% тул Gemini-г сонгосон гэж бич; API-г нь бич»): 11/15-нд унтраасан 12 бичлэгийн туршилтыг
% full/ch3_evaluation_and_stt_choice.tex-ээс энд шилжүүлж, API-ийн тайлбарыг нэмэв (gemini_stt_runtime.py,
% tools/compare_stt_candidate.py, demo/index.html-тэй тулгасан). Үр дүн нь 4.2-т. Мөн багшийн саналаар
% «Бүлгийн дүгнэлт»-ийг хасав (ШМТК-ийн дипломд байдаггүй).
